\documentclass[10pt,a4paper]{amsart}
\usepackage[margin=19mm]{geometry}
\usepackage{amsmath,amssymb,mathtools}
\usepackage[scr=rsfs,cal=euler]{mathalfa}
\usepackage{xfrac}
\usepackage[unicode,hidelinks,bookmarks=false]{hyperref}
\newcommand{\ie}{i.e.\ }
\newcommand{\cf}{cf.\ }
\newcommand{\resp}{respectively\ }
\newcommand{\Subsection}{\subsection}
\providecommand{\nobreakdash}{\nobreak\hbox{-}}
\setlength{\emergencystretch}{2em}
\multlinegap0pt
\usepackage{amssymb}
\usepackage{xcolor}
\usepackage{tikz}
\usepackage{mathtools}
\newtheorem{prop}{Proposition}
\newcommand{\as}{\boldsymbol{\alpha}}
\newcommand{\bR}{\mathbb{R}}
\newcommand{\bb}{\boldsymbol{b}}
\newcommand{\bs}{\boldsymbol{\beta}}
\newcommand{\cH}{\mathcal{H}}
\newcommand{\cI}{\mathcal{I}}
\newcommand{\one}{\boldsymbol{1}}
\newcommand{\qs}{\boldsymbol{q}}
\newcommand{\xs}{\boldsymbol{x}}
\newcommand{\ys}{\boldsymbol{y}}
\title{Quantum braid Floer homology}
\author{Cristina Ana-Maria Anghel, Andr\'as Juh\'asz}
\date{Source: arXiv:2601.08805v4 (CC BY 4.0)}
\begin{document}
\maketitle

\noindent\emph{Source and licence.} Quantum braid Floer homology. Version arXiv:2601.08805v4 (CC BY 4.0), distributed by arXiv under the Creative Commons Attribution 4.0 International licence, http://creativecommons.org/licenses/by/4.0/, as recorded in the arXiv licence metadata for this exact version and shown on the abstract page https://arxiv.org/abs/2601.08805v4. Full licence notice: \texttt{licenses/arxiv-2601.08805-CC-BY-4.0.txt}.\par\smallskip

\noindent\textit{Editorial scope.} Counted unit: the article's prop prop:W-S (source0.tex lines 1561--1567). The article's complete original proof of it is delivered with it (source0.tex lines 1569--1611, one \texttt{proof} environment of 43 lines). No mathematics is written, repaired or re-derived: the statement and the proof are the article's own lines, in the article's own notation, and no retained line is substituted or re-flowed.  No further source-local context is needed: every cross-reference of every retained line resolves inside the two delivered spans.  Nothing was omitted from any retained span, so the package carries no omission record.  No retained line contains a citation, so the wrapper carries no bibliography and no bibliography entry.  No retained line contains a figure environment, an image-inclusion command or drawing code, so no image is delivered and the package declares no asset dependency.  Editorial formatting only, in addition: an AMS \texttt{amsart} preamble that restates the article's own declarations and macros used by the retained lines, namely the environments \texttt{prop} on separate counters, together with the standard packages those lines need.  The retained environments print in this document as Proposition 1 in order of appearance, whereas the article prints the same items with the numbering of its own full document.  The complete native source of the article as distributed in the arXiv e-print for this exact version is bundled unchanged as \texttt{sources/192630/source0.tex}.
\begin{prop}\label{prop:W-S}
    Let $\beta \in B_{n+1}$ be a braid whose closure is a knot and $\cH^q_\beta$ an associated quantum Heegaard diagram. Suppose that $\xs$, $\ys \in \cI$, and $D \in D_{\cH}(\xs,\ys)$ is such that $n_w(D) = 0$. Then
    \begin{equation}\label{eqn:W-S}
    W(\ys) - W(\xs) = 3n_{\qs}(D) + |S(\ys)| - |S(\xs)|.
    \end{equation}
    Furthermore, if $D \ge 0$, then $|S(\xs)| \le |S(\ys)|$ and $W(\xs) \le W(\ys)$.
\end{prop}

\begin{proof}
    Given $H \subseteq \{1,\dots,n\}$, we write $\one_H$ for the indicator function of $H$. Set
    \[
    \omega_i(D) := \one_{S(\ys)}(i) - \one_{S(\xs)}(i) \in \{-1,0,1\}.
    \]
    Then
    \begin{equation}\label{eqn:omega-S}
    \sum_{i=1}^n \omega_i(D) = |S(\ys)| - |S(\xs)|.
    \end{equation}
    For $i \in \{1,\dots,n\}$, let $m_i(D)$ be the multiplicity of 
    \[
    a_i \cap \left\{x \in \bR^2 : 0 < x < \frac{1}{2(n+2)}\right\} \subseteq \alpha_i
    \]
    in $\partial_\alpha D := \partial D \cap \as$. Since $a_i \cap \bb \cap \{x>0\} = \{x_i^0\}$, the equation $\partial(\partial_\alpha D) = \ys - \xs$ gives $m_i(D) = \omega_i(D)$.

    As before, let $q_{n+1}$ be a point in the region neighbouring the one containing $q_n$ along $a_n$. For $i \in \{1,\dots,n+1\}$, join $w$ to $q_i$ by a path $\nu_i$ such that $\nu_i \pitchfork (\as \cup \bs)$ consists of $i-1$ points lying successively in $a_1,\dots,a_{i-1}$. Since $n_w(D) = 0$, the change in the coefficient of $D$ along $\nu_i$ gives
    \begin{equation}\label{eqn:omega-nq-i}
    n_{q_i}(D) = \sum_{j=1}^{i-1} \omega_j(D).
    \end{equation}
    As $q_0$ and $w$ lie in the same region, $n_{q_0}(D) = 0$. We obtain that
    \begin{equation}\label{eqn:omega-nq}
    n_{\qs}(D) = \sum_{i=1}^n n_{q_i}(D) = \sum_{i=1}^n \sum_{j=1}^{i-1} \omega_j(D) = \sum_{j=1}^n (n-j)\omega_j(D).
    \end{equation}
    By definition,
    \[
    W_i(\xs) = 
    \begin{cases}
        3(n-i)+1, & x_i = x_i^0,\\
        0, & x_i \neq x_i^0.
    \end{cases}
    \]
    Hence,
    \[
    W(\ys) - W(\xs) = \sum_{i=1}^n (3(n-i)+1) \omega_i(D) = 3\sum_{i=1}^n (n-i)\omega_i(D) + \sum_{i=1}^n \omega_i(D).
    \]
    Using equations~\eqref{eqn:omega-S} and~\eqref{eqn:omega-nq}, this becomes equation~\eqref{eqn:W-S}.

    Now assume that $D \ge 0$. By equation~\eqref{eqn:omega-nq-i}, we have 
    \[
    0 \le n_{q_{n+1}}(D) = \sum_{j=1}^n \omega_j(D) = |S(\ys)| - |S(\xs)|. 
    \]
    This, together with $n_{\qs}(D) \ge 0$ and equation~\eqref{eqn:W-S} implies that $W(\ys) - W(\xs) \ge 0$.
\end{proof}

\end{document}
